\chapter{Controle Sintético}
\label{cap:synth}
% ── índice ──────────────────────────────────────────────────────
\index{controle sintético|textbf}
\index{synth@\texttt{synth()}|textbf}
\index{RMSPE|textbf}
\index{doadores, pool de}
\index{placebo, testes de}

% ─────────────────────────────────────────────────────────────────────────────
\section{O problema do caso único}

O DiD (cap.~\ref{cap:did}) e o PSM (cap.~\ref{cap:psm}) dependem de múltiplas
unidades tratadas e controles comparáveis. Em muitas aplicações de política,
porém, apenas uma unidade recebe o tratamento — um país que adota um programa
econômico, um estado que muda sua legislação — e as unidades potencialmente
comparáveis são poucas e heterogêneas.

O \textbf{Controle Sintético} (Abadie, Diamond e Hainmueller, ADH 2010)
constrói um grupo de controle artificial — o \emph{controle sintético} — como
uma combinação convexa de unidades doadoras que replica a trajetória da unidade
tratada \emph{antes} do tratamento. A ideia central:

\[
  \hat{Y}_{1t}^{(0)} = \sum_{j=2}^{J+1} w_j^*\, Y_{jt}
  \qquad
  w_j^* \geq 0,\quad \sum_{j=2}^{J+1} w_j^* = 1
\]

Os pesos $W^* = (w_2^*,\ldots,w_{J+1}^*)$ minimizam o erro quadrático médio
de predição (\emph{RMSPE}) no período pré-tratamento:

\[
  W^* = \arg\min_{W \in \Delta^J}
  \sum_{t=1}^{T_0-1} \left( Y_{1t} - \sum_{j\geq 2} w_j Y_{jt} \right)^2
\]

onde $\Delta^J$ é o simplexo unitário ($w_j \geq 0$, $\sum w_j = 1$). O efeito
do tratamento no período $t \geq T_0$ é:

\[
  \hat\tau_t = Y_{1t} - \hat{Y}_{1t}^{(0)}
\]

A validade do método repousa sobre a qualidade do ajuste pré-tratamento:
se o controle sintético acompanha a unidade tratada no período pré-tratamento,
a hipótese de que teria continuado a fazê-lo na ausência do tratamento é
plausível — uma forma não paramétrica da hipótese de tendências paralelas.

\begin{nota}
  A restrição de convexidade ($w_j \geq 0$, $\sum w_j = 1$) evita
  extrapolação — o controle sintético fica dentro do suporte dos doadores.
  Isso o distingue do DiD com pesos, que pode atribuir pesos negativos para
  alcançar o balanço.
\end{nota}

% ─────────────────────────────────────────────────────────────────────────────
\section{Verificação por DGP}

O DGP usa painel de 10 unidades $\times$ 20 períodos. A unidade 1 é tratada
a partir do ano 11; as unidades 2--10 são doadoras. O efeito verdadeiro é
$\tau = 3{,}0$ (permanente, a partir de $T_0 = 11$).

\[
  y_{it} = \alpha_i + 0{,}3\,t + 3{,}0\cdot\mathbf{1}[i=1, t\geq 11]
           + \varepsilon_{it}
  \qquad
  \alpha_i \in \{2,3,\ldots,10\},\quad \alpha_1 = 5
\]

O intercepto da unidade tratada ($\alpha_1 = 5$) coincide com o da unidade 5
(doadora), garantindo suporte comum. As inovações $\varepsilon_{it} \sim U(0,1)$
introduzem ruído que impede o SC de usar apenas a unidade 5 como doadora
perfeita.

\begin{lstlisting}[caption={DGP: painel 10 × 20, tratamento na unidade 1 a partir de t=11}]
seed(42)
input df
id
1
end
for i in 1..=8 { df = append(df, df) }
df |> filter(_n <= 200)

generate df unit = (_n - 1) % 10 + 1
generate df year = (_n - (_n - 1) % 10 - 1) / 10 + 1

// interceptos por unidade: doadores recebem alpha = unit; tratado alpha = 5
generate df alpha = 0.0
let j = 2
while j <= 10 {
    replace df alpha = j * 1.0 if unit == j
    j = j + 1
}
replace df alpha = 5.0 if unit == 1

generate df d = (unit == 1) * (year >= 11)
generate df e = rnormal()
generate df y = alpha + 0.3 * year + 3.0 * d + e

let m = synth("y", 1, 11, df, id="unit", time="year")
display m
\end{lstlisting}

\subsection*{Construção do painel em Hayashi}

A função \lstinline{generate} opera sobre o vetor \lstinline{_n} (linha
corrente, 1-indexado). Para um painel $N \times T$ em formato longo
(empilhado por unidade), as fórmulas são:

\begin{center}
\begin{tabular}{ll}
\toprule
\textbf{Variável} & \textbf{Expressão} \\
\midrule
Unidade ($1,\ldots,N$)    & \lstinline{(_n - 1) % N + 1} \\
Período ($1,\ldots,T$)    & \lstinline{(_n - (_n-1)%N - 1) / N + 1} \\
\bottomrule
\end{tabular}
\end{center}

\subsection*{Resultado}

\begin{lstlisting}[language={}, basicstyle=\ttfamily\small,
  frame=none, numbers=none, backgroundcolor=\color{codbg},
  xleftmargin=0pt, xrightmargin=0pt, breaklines=false]
════════════════════════════════════════════════════════════════════
 Controle Sintético  —  Abadie, Diamond, Hainmueller (2010)
════════════════════════════════════════════════════════════════════
 Unidade tratada : 1   Outcome: y
 T₀ (1ª pós-trat): 11   Doadores: 9   T pré: 10   T pós: 10
 RMSPE pré  : 0.3347
 RMSPE pós  : 2.7056   razão pós/pré: 8.083
────────────────────────────────────────────────────────────────────
 Pesos dos doadores (w > 0.001):
   6                         0.4420
   3                         0.2153
   8                         0.1747
   2                         0.1406
   7                         0.0274
────────────────────────────────────────────────────────────────────
  Período          Real     Sintético        Efeito
 ──────────────────────────────────────────────────
         1        5.8266        6.0335
         2        6.5321        6.1112
         3        6.7458        6.5023
         4        6.6240        6.8732
         5        7.4867        7.2037
         6        7.5689        7.4293
         7        8.0670        7.6433
         8        7.8780        8.0722
         9        8.0202        8.3241
        10        8.0878        8.6960
 *      11       11.4437        8.8167        2.6270
 *      12       12.1888        9.4361        2.7527
 *      13       12.0728        9.8984        2.1745
 *      14       12.9748        9.5753        3.3995
 *      15       12.7653       10.2134        2.5519
 *      16       13.2120       10.2424        2.9696
 *      17       13.3167       10.6701        2.6465
 *      18       13.7205       11.0759        2.6446
 *      19       14.3320       11.5282        2.8039
 *      20       14.0020       11.7106        2.2913
════════════════════════════════════════════════════════════════════
 * pós-tratamento
\end{lstlisting}

\subsection*{Leitura dos resultados}

\paragraph{Ajuste pré-tratamento.}
O RMSPE pré-tratamento é $0{,}33$ em uma série cujo desvio padrão é da
ordem de 2.5 — um ajuste razoável dado o ruído do DGP ($N=200$, apenas
10 doadores). A razão pós/pré = 8.1 sinaliza que a divergência
pós-tratamento é \emph{oito vezes} maior do que o erro de ajuste
pré-tratamento, rejeitando a hipótese de efeito nulo.

\paragraph{Pesos.}
O algoritmo distribui os pesos principalmente entre as unidades 6, 3, 8 e
2, em vez de concentrar todo o peso na unidade 5 (que tem $\alpha_5 = 5 =
\alpha_1$). Isso ocorre porque o ruído amostral $\varepsilon_{it}$ torna a
trajetória realizada da unidade 5 uma combinação linear imprecisa da
unidade 1 no período pré-tratamento. O SC usa múltiplos doadores para
cancelar o ruído.

\paragraph{Efeitos estimados.}
A média dos efeitos pós-tratamento é:

\[
  \bar{\hat\tau} = \frac{1}{10}\sum_{t=11}^{20}\hat\tau_t = 2{,}69
\]

O valor verdadeiro é $\tau = 3{,}0$; a subestimação de ${\sim}10\%$ reflete
a imprecisão do ajuste pré-tratamento com apenas 10 doadores e 10 períodos
pré.

% ─────────────────────────────────────────────────────────────────────────────
\section{Sintaxe}

\begin{lstlisting}
synth("outcome", treated_id, t0, df, id="col_unidade", time="col_tempo")
synth("outcome", treated_id, t0, df, id="col", time="col", covs=["x1","x2"])
\end{lstlisting}

\begin{center}
\begin{tabular}{lll}
\toprule
\textbf{Argumento} & \textbf{Tipo} & \textbf{Descrição} \\
\midrule
\texttt{outcome}    & string & nome da variável de resultado \\
\texttt{treated\_id}& num/str& identificador da unidade tratada \\
\texttt{t0}         & número & primeiro período pós-tratamento \\
\texttt{df}         & DataFrame & painel em formato longo \\
\texttt{id=}        & opção  & coluna de identificação de unidade \\
\texttt{time=}      & opção  & coluna de período \\
\texttt{covs=}      & lista  & covariáveis adicionais para o ajuste \\
\bottomrule
\end{tabular}
\end{center}

O argumento \texttt{covs} permite incluir médias de covariáveis pré-tratamento
no critério de minimização, além dos valores da variável de resultado. Útil
quando os doadores têm trajetórias similares mas características estruturais
distintas.

\begin{lstlisting}[caption={Fluxo típico — avaliação de política estadual}]
load "estados.csv" as df   // id, ano, pib, pop, educ

let m = synth("pib", "SP", 2008, df,
              id="id", time="ano",
              covs=["pop", "educ"])
display m
\end{lstlisting}

\begin{nota}
  O CS não fornece inferência assintótica padrão (sem erros-padrão, sem
  $p$-values). A inferência é feita via \emph{testes de placebo}: aplica-se
  o SC a cada unidade doadora "como se" fosse tratada e compara a razão
  pós/pré. Se a razão para a unidade verdadeiramente tratada for excepcional
  no conjunto de placebos, o efeito é estatisticamente incomum.
  O \hayashi{} exibe a razão pós/pré como heurística descritiva; testes de
  placebo são implementados manualmente via loop sobre as unidades.
\end{nota}

\section{Validação empírica}

O caso de validação \texttt{synth\_smoking} em \texttt{validation/cases/}
estima um controle sintético sobre um painel simulado (10 doadores, 1 unidade
 tratada, 20 períodos). O caso compara o efeito pós-tratamento médio do
\hayashi{} com uma implementação de referência em Python. Execute
\lstinline{hay validate} para ver o estado atual.
